\documentclass[11pt,a4paper]{ctexart}

% ==================================================
% 1. 宏包引入
% ==================================================
\usepackage{hyperref}
\usepackage{xcolor}
\usepackage{calc}
\usepackage{graphicx}
\usepackage{tikz}
\usepackage{fontspec}
\usepackage{fontawesome5}
\usepackage{titlesec}
\usepackage{enumitem}
\usepackage{fancybox}
\usepackage{geometry}
\usepackage{tabularx} % <--- 新增：引入 tabularx 宏包

% 引入 TikZ 计算库
\usetikzlibrary{calc}

\hypersetup{hidelinks}

% ==================================================
% 2. 全局设置与格式
% ==================================================
\setlength{\parindent}{0pt}					% 取消全局段落缩进
\pagenumbering{gobble}						% 取消页码显示
\setlist[itemize]{nosep                     % 取消 itemize 的默认间距
    , before={\vspace*{-\parskip}}          % 取消 itemize 和后续段落之间的空白
    , leftmargin=*}		                    % 取消 itemize 的左边距
\setlist[enumerate]{leftmargin=*}	        % 取消 enumerate 的左边距
\renewcommand{\arraystretch}{1.2}           % 设置表格行间距
\linespread{1.25}                           % 设置正文行间距

% 标题格式设置（下方带横线）
\titleformat{\section}
  {\LARGE\bfseries\raggedright} 		      
  {}{0em}                      			  
  {}                           			  
  [{\color{secondary_color}\titlerule}]     
\titlespacing*{\section}{0cm}{*1.2}{*1.2}	

% 页面边距设置（完全参照东南大学排版）
\geometry{
	a4paper,
	left=1.2cm,
	right=1.2cm,
	top=1.5cm,
	bottom=1cm,
	nohead
}

% 字体设置（如果报错找不到字体，可以注释掉这段使用默认字体）
\setmainfont[
    Path=fonts/,
    Extension=.otf,
    BoldFont=*-Bold,
]{NotoSerifSC}

% ==================================================
% 3. 颜色定制（东北大学深蓝色）
% ==================================================
\definecolor{neublue}{HTML}{1E3C96}    
\definecolor{primary_color}{HTML}{1E3C96}    % 主色调：深蓝
\definecolor{secondary_color}{HTML}{1E3C96}  % 副色调（标题下划线）：深蓝

\newlength{\iconwidth}
\setlength{\iconwidth}{1.5em}                   % 设置 section 标题部分图标占用的宽度

% ==================================================
% 4. 图片命名与路径（名称不可变）
% ==================================================
\graphicspath{{images/}}
\newcommand{\HeaderLogoFile}{neu-header.png}
\newcommand{\WatermarkLogoFile}{neu-watermark.png}
\newcommand{\PhotoFile}{photo.jpg}

% ==================================================
% 5. 页眉、页脚与背景（通过 Hook 实现深蓝背景与精准定位）
% ==================================================
\AddToHook{shipout/background}{%
\begin{tikzpicture}[remember picture, overlay]
  \coordinate (NW) at (current page.north west);
  \coordinate (NE) at (current page.north east);
  \coordinate (SW) at (current page.south west);
  \coordinate (SE) at (current page.south east);

  % 顶部深蓝色背景 (高度1.75cm)
  \fill[neublue] (NW) rectangle ([yshift=-1.75cm]NE);

  % 底部深蓝色背景 (高度0.65cm)
  \fill[neublue] (SW) rectangle ([yshift=0.65cm]SE);

  % 背景水印（完美居中，透明度10%）
  \node[anchor=center, inner sep=0pt, opacity=0.10] at (current page.center) {%
    \includegraphics[width=18cm, height=18cm, keepaspectratio]{\WatermarkLogoFile}%
  };
\end{tikzpicture}%
}

\AddToHook{shipout/foreground}{%
\begin{tikzpicture}[remember picture, overlay]
  \coordinate (NW) at (current page.north west);
  \coordinate (NE) at (current page.north east);
  \coordinate (S)  at (current page.south);

  % 左侧：东北大学校徽与校名
  \node[anchor=west, inner sep=0pt] at ([xshift=1.55cm, yshift=-0.88cm]NW) {%
    \includegraphics[width=5.25cm, height=1.30cm, keepaspectratio]{\HeaderLogoFile}%
  };

  % 右侧：学院名称
  \node[anchor=east, text=white, font=\bfseries\small] at ([xshift=-1.55cm, yshift=-0.76cm]NE) {悉尼智能科技学院};
  \node[anchor=east, text=white, font=\footnotesize] at ([xshift=-1.55cm, yshift=-1.17cm]NE) {Sydney Smart Technology College / SSTC};

  % 底部：联系方式
  \node[text=white, font=\footnotesize] at ([yshift=0.32cm]S) {%
    \faPhone\ 130-6666-0000
    \qquad
    \faEnvelope\
    \href{mailto:youremail@neu.edu.cn}{youremail@neu.edu.cn}
    \qquad
    \faGithub\
    \href{https://github.com/你的GitHub}{GitHub 项目地址}
  };
\end{tikzpicture}%
}

% ==================================================
% 6. 文档正文
% ==================================================
\begin{document}

    % -------------------- 个人信息与照片 --------------------
    \begin{minipage}[t]{0.78\textwidth}
        \begin{minipage}[t]{\textwidth}
        \section[个人信息]{\makebox[\iconwidth][c]{\color{primary_color}{\faAddressCard}}\quad 个人信息}
        
        \begin{minipage}[t]{0.5\textwidth}
            \textbf{姓\qquad 名}：你的名字
            
            \vspace{0.5em}
            \textbf{出生年月}：你的出生日期
        \end{minipage}
        \begin{minipage}[t]{0.35\textwidth}
            \textbf{性\qquad 别}：沃尔玛购物袋
            
            \vspace{0.5em}
            \textbf{政治面貌}：你的政治面貌
        \end{minipage}
        \vspace{1.2em}
        \end{minipage}

        % -------------------- 教育背景 --------------------
        \begin{minipage}[t]{\textwidth}
        \section[教育背景]{\makebox[\iconwidth][c]{\color{primary_color}{\faGraduationCap}}\quad 教育背景}
        
        {\large \textbf{东北大学}}，专科 \hfill 2000年9月--2010年6月
        \begin{itemize}
            \item 你的学院，你的专业
            \item \textbf{主修课程}：课程1、课程2、课程3、课程4\ 等。
        \end{itemize}
        
        \vspace{0.5em}
        {\large \textbf{东北大学}}，本科 \hfill 2010年9月--2020年6月
        \begin{itemize}
            \item 你的学院，你的专业
            \item \textbf{主修课程}：课程1、课程2、课程3、课程4\ 等。
            \item \textbf{GPA}：4.8 / 4.8（排名：1 / 250）
        \end{itemize}
        
        \vspace{0.5em}
        {\large \textbf{东北大学}}，硕士 \hfill 2020年9月--至今
        \begin{itemize}
            \item 你的学院，你的专业，你的导师姓名\ 职称
            \item \textbf{研究方向}：方向1、方向2、方向3、方向4\ 等。
        \end{itemize}
        
        \vspace{1.2em}
        \end{minipage}
    \end{minipage}
    \hfill
    % 右半边照片（带占位符兜底机制）
    \begin{minipage}[t]{0.2\textwidth}
        \vspace{2em} 
        \setlength{\fboxsep}{0pt}
        \IfFileExists{images/\PhotoFile}{
            \doublebox{\includegraphics[width=\linewidth]{\PhotoFile}}
        }{
            \doublebox{\begin{tikzpicture}
                \draw[neublue!55, dashed] (0,0) rectangle (\linewidth, 3.5cm);
                \node[align=center, text=gray, font=\small] at (0.5\linewidth, 1.75cm) {证件照\\photo.jpg};
            \end{tikzpicture}}
        }
    \end{minipage}

    % -------------------- 科研成果 --------------------
    \begin{minipage}[t]{\textwidth}
    \section[科研成果]{\makebox[\iconwidth][c]{\color{primary_color}{\faAtom}}\quad 科研成果}

    This is One of Your Paper Published in Conference A.
    \begin{itemize}
        \item \textbf{Mingzi Nide}, Daoshi Nide. \hfill 发表于 \textbf{Conference A}（CCF-A类会议）
        \item 用一句话描述这份论文干了什么\dots\dots
    \end{itemize}

    \vspace{0.5em}
    《另一份论文的标题》
    \begin{itemize}
        \item  \textbf{你的名字}、你的师兄、你的导师 \hfill 发表于 \textbf{某篇期刊} （SCI-1区）
        \item 用一句话描述这份论文干了什么\dots\dots
    \end{itemize}
    
    \vspace{1.2em}
    \end{minipage}

    % -------------------- 项目与教学 --------------------
    \begin{minipage}[t]{\textwidth}
    \section[项目与教学]{\makebox[\iconwidth][c]{\color{primary_color}{\faChalkboardTeacher}}\quad 项目与教学}
    
    {\large \textbf{项目名称}} \hfill 2020年9月--2021年9月
    \begin{itemize}
        \item \textbf{你在项目中扮演的角色} \hfill 横向/纵向项目-已完结/进行中
        \item 用一句话介绍你在这个项目中做了什么\dots\dots
    \end{itemize}

    \vspace{0.5em}
    {\large \textbf{某某主题讨论班}}，主讲 / 参与 \hfill 2020年夏季
    \begin{itemize}
        \item \textbf{主要内容}：内容1，内容2，内容3\ 等。
    \end{itemize}

    \vspace{0.5em}
    {\large \textbf{课程名称}}，助教 \hfill 2021年夏季
    \begin{itemize}
        \item \textbf{主要内容}：内容1，内容2，内容3\ 等。
    \end{itemize}
    
    \vspace{1.2em}
    \end{minipage}
    
    % -------------------- 技能特长 & 兴趣爱好 --------------------
    \begin{minipage}[t]{0.6\textwidth}
        \section[技能特长]{\makebox[\iconwidth][c]{\color{primary_color}{\faWrench}}\quad 技能特长}
        \begin{itemize}
        \setlength{\itemsep}{0.5em}
            \item 熟练使用 Python、Jvav、Rust 等编程语言。
            \item 熟练使用 Tensorflow、Pytorch 等深度学习框架。
            \item 熟悉 Windows 与 Linux 端开发。
        \end{itemize}
    \end{minipage}
    \hfill
    \begin{minipage}[t]{0.35\textwidth}
        \section[兴趣爱好]{\makebox[\iconwidth][c]{\color{primary_color}{\faStar}}\quad 兴趣爱好}
        \begin{itemize}
        \setlength{\itemsep}{0.5em}
            \item 我喜欢唱
            \item 跳
            \item Rap
            \item 篮球
        \end{itemize}
    \end{minipage}

\newpage
    % % 如有需要，可以添加额外的页面。不要忘记添加页眉页脚和背景相关的代码。

    % % -------------------- 竞赛经历 --------------------
    % \begin{minipage}[t]{\textwidth}
    % \section[竞赛经历]{\makebox[\iconwidth][c]{\color{primary_color}{\faTrophy}}\quad 竞赛经历}
    % \vspace{0.5em}
    
    % % 使用 tabularx 环境，X 列会自动拉伸并填满剩余宽度
    % \begin{tabularx}{\textwidth}{X c c c}
    %     \textbf{比赛1} & 第一负责人 & 国家级-第10名 & 2023年4月 \\
    %     \textbf{比赛2} & 个人参赛 & 国家级-一等奖 & 2023年8月 \\
    %     \textbf{比赛3} & 个人参赛 & 省级-一等奖 & 2022年12月 \\
    % \end{tabularx}
    
    % \vspace{1.2em}
    % \end{minipage}

    % % -------------------- 技能特长 --------------------
    % \begin{minipage}[t]{\textwidth}
    % \section[技能特长]{\makebox[\iconwidth][c]{\color{primary_color}{\faWrench}}\quad 技能特长}
    % \begin{itemize}
    %     \item 熟练使用 C++、Python、Matlab 编程语言。
    %     \item 熟悉 Windows 与 Linux 端开发。
    %     \item 熟练使用 Tensorflow，Pytorch 等深度学习框架。
    %     \item 熟练掌握 C++ 与 Python 环境下 OpenCV 与 Qt 应用的开发，且熟练使用 Qt Creator 软件。
    %     \item 熟练使用 Altium Designer 与 LCEDA 进行封装绘制与板子设计。
    %     \item 熟练使用 Keil，Arduino IDE 等集成开发软件。
    %     \item 了解模式识别，强化学习，遗传算法，知识蒸馏等相关概念。
    % \end{itemize}
    % \vspace{1.2em}
    % \end{minipage}

    % % -------------------- 所获荣誉 --------------------
    % \begin{minipage}[t]{\textwidth}
    % \section[所获荣誉]{\makebox[\iconwidth][c]{\color{primary_color}{\faStar}}\quad 所获荣誉}
    % \vspace{0.5em}
    
    % % 使用两个并排的 minipage 代替 multicols 环境
    % \begin{minipage}[t]{0.48\textwidth}
    %     \begin{itemize}
    %         \item 某年学业先进个人
    %         \item 某年某奖学金某等奖
    %         \item 某大使
    %     \end{itemize}
    % \end{minipage}
    % \hfill
    % \begin{minipage}[t]{0.48\textwidth}
    %     \begin{itemize}
    %         \item 某年某奖学金某等奖
    %         \item 某年优秀团员称号
    %         \item 某年某称号
    %     \end{itemize}
    % \end{minipage}
    
    % \vspace{1.2em}
    % \end{minipage}

    % % -------------------- 其他 --------------------
    % \begin{minipage}[t]{\textwidth}
    % \section[其他]{\makebox[\iconwidth][c]{\color{primary_color}{\faInfo}}\quad 其他}
    % \begin{itemize}
    %     \item 英语水平-CET6级xxx分
    %     \item 计算机几级证书
    %     \item xx几级证书
    %     \item 技术博客: 某网址
    %     \item 教师资格证：xxx
    %     \item 普通话证书：几级几等
    %     \item 文字排版：\LaTeX
    % \end{itemize}
    % \end{minipage}

\end{document}